\documentclass{article}
\usepackage{/home/user1/code/tex/sty}
\usepackage{amsfonts}
\usepackage{amsmath}
\usepackage{mathtools}
\usepackage{geometry}
\usepackage{graphicx}
\usepackage{hyperref}
\usepackage{floatrow}
\newcommand{\fig}[3]{
\begin{figure}[H]
\begin{center}
\includegraphics[height=#1]{#2}
\caption{#3}
\label{#2}
\end{center}
\end{figure}}
\setlength{\parindent}{0pt}
\geometry{top=1in,bottom=1in,left=1in,right=1in}
\begin{document}
\begin{center}
\textbf{\Large{Approximate Solutions of the Navier Stokes Equations}}\\~\\
\begin{tabular}{l l}
Student&Agustin Romero\\
Email&ar1@stanford.edu\\
Instructor&Lernik Asserain\\
Course&MATH 131P\\
Institution&Stanford University\\
Date&\today\\
\end{tabular}
\end{center}
\begin{abstract}
\end{abstract}
The Navier Stokes equations with radial symmetry are derived and solution found as a Bessel-Fourier series.
\section{Navier Stokes Equations in Rectangular Coordinates}
Consider the three dimensional box enclosing a fluid. The position of and particle in the fluid is given by the solution $u(t,x,y,z)$, whose velocity field is $\vec{V}(x,y,z)$.
\fig{5cm}{1.png}{Three dimensional box with all stresses in the x direction shown to first order.}
Assume there are only two forces on the fluid: 1) the force of gravity, 2) the force from the fluid along the boundaries. The force on the fluid by gravity along the $a$ axis is
\e{\rho g_a}
where $\rho$ is the mass density of the volume element. The force on the fluid from the fluid is quantified by the shear stress $\tau_{ab}$: the rate of change in position of the fluid in the $b$ axis, given a force acting in the $y$ axis. Suppose that, for each face of the box, the shear stress $\tau$ is approximated by a Taylor expansion to linear order, so that the shear stress is at most linearly dependent on position. The total force on the fluid along the $a$ axis $dF_a$ the sum of these two contributions. 
\e{dF_x=(\rho g_x + \pa{x}\tau_{xx}+\pa{y}\tau_{yx}+\pa{z}\tau_{zx})dxdydz}
\e{dF_y=(\rho g_y + \pa{x}\tau_{xy}+\pa{y}\tau_{yy}+\pa{z}\tau_{zy})dxdydz}
\e{dF_z=(\rho g_z + \pa{x}\tau_{xz}+\pa{y}\tau_{yz}+\pa{z}\tau_{zz})dxdydz}
This expresses the net force on the particle in all directions. We are now in a position to apply Newton's force equation $\vec{F}=m\vec{a}$. Newton's second law is
\e{\vec{F}&=m\vec{a}\\
&=\dd{t}\vec{P}\\
&= dm (u\pa{x} \vec{V} + v\pa{y} \vec{V} + w\pa{z} \vec{V} + \pa{t}\vec{V})}
where $\vec{P}$ is the linear momentum of the system.
\e{\vec{P}=\int_M \vec{V} dm}
where $dm$ is the infinitesimal fluid particle of mass $dm$, and $V$ is the particle velcoity in the fluid.
\e{d\vec{F}=dm \dd{t}\vec{V}}
We need a way to quantify the acceleration field $\vec{a}$. Use the time derivative on $\vec{V}$ and apply the chain rule.
\e{\vec{a}&=\dd{t}\vec{V}=\pa{t}\vec{V}+\pa{x}\vec{V}\dd{t}x+\pa{y}\vec{V}\dd{t}y+\pa{z}\vec{V}\dd{t}z\\
&=\pa{t}\vec{V}+u\pa{x}\vec{V}+v\pa{y}\vec{V}+w\pa{z}\vec{V}\\
&=\pa{t}\vec{V}+(u\pa{x}+v\pa{y}+w\pa{z})\vec{V}}
This is intuitively explained by two types of contributions: 1) the change in the velocity field $\vec{V}$ per unit time; 2) the change in position of the velocity field $\vec{V}$ by deformation. Now apply $\vec{F}=m\vec{a}$ infinitesimally.
\e{d\vec{F}&=\rho d\vec{a}\\
\rho g_x + \pa{x}\tau_{xx}+\pa{y}\tau_{yx}+\pa{z}\tau_{zx}&=\rho(\pa{t}u+u\pa{x}u+v\pa{y}u+w\pa{z}u)\\
\rho g_y + \pa{x}\tau_{xy}+\pa{y}\tau_{yy}+\pa{z}\tau_{zy}&=\rho(\pa{t}v+u\pa{x}v+v\pa{y}v+w\pa{z}v)\\
\rho g_z + \pa{x}\tau_{xz}+\pa{y}\tau_{yz}+\pa{z}\tau_{zz}&=\rho(\pa{t}w+u\pa{x}w+v\pa{y}w+w\pa{z}w)}

This is the differential momentum equation.
\e{\rho \fr{D}{Dt}u = \rho g_x - \pa{x} p + \pa{x}(\mu (2\pa{x} u - \fr{2}{3}\nabla \cdot \vec{V})) + \pa{y}(\mu (\pa{y} u + \pa{x} v)) + \pa{z}(\mu (\pa{x} w + \pa{z} u))}
\e{\rho \fr{D}{Dt}v = \rho g_y - \pa{y} p + \pa{x}(\mu (\pa{y} u + \pa{x} v)) + \pa{y}(\mu (2\pa{y} v - \fr{2}{3}\nabla \cdot \vec{V})) + \pa{z}(\mu (\pa{z} v + \pa{y} w))}
\e{\rho \fr{D}{Dt}w = \rho g_z - \pa{z} p + \pa{x}(\mu (\pa{x} w + \pa{z} u)) + \pa{y}(\mu (\pa{z}w+\pa{y}w)) + \pa{z}(\mu (2\pa{z} w - \fr{2}{3}\nabla \cdot \vec{V}))}
These are the Navier Stokes equations. They assume the fluid is influced by linear shear stresses, linear deformations, and gravity, coupled by Newton's force equation $\vec{F}=m\vec{a}$. For incompressible flow,
\e{\nabla \cdot \vec{V}=0}
For non-convective flow,
\e{\fr{D}{Dt}\vec{V}=(\vec{V}\cdot \nabla)\vec{V}+\pa{t}\vec{V}=\pa{t}\vec{V}}
Assuming incompressible and non-convective flow, the Navier-Stokes equations are,
\e{\rho(\pa{t}u + u\pa{x}u + v\pa{y}u + w\pa{z}u)=\rho g_x - \pa{x}p + \mu(\ppa{x}{2}u + \ppa{y}{2}u + \ppa{z}{2}u)}
\e{\rho(\pa{t}v + u\pa{x}v + v\pa{y}v + w\pa{z}v)=\rho g_y - \pa{y}p + \mu(\ppa{x}{2}v + \ppa{y}{2}v + \ppa{z}{2}v)}
\e{\rho(\pa{t}w + u\pa{x}w + w\pa{y}w + w\pa{z}w)=\rho g_z - \pa{z}p + \mu(\ppa{x}{2}w + \ppa{y}{2}w + \ppa{z}{2}w)}
For frictionless flow,
\e{\mu=0}
Assuming incompressible, constant-viscosity, non-convective, and frictionless flow, Navier-Stokes equations are
\e{\rho(\pa{t}u + u\pa{x}u + v\pa{y}u + w\pa{z}u)=\rho g_x - \pa{x}p}
\e{\rho(\pa{t}v + u\pa{x}v + v\pa{y}v + w\pa{z}v)=\rho g_y - \pa{y}p}
\e{\rho(\pa{t}w + u\pa{x}w + w\pa{y}w + w\pa{z}w)=\rho g_z - \pa{z}p}

\section{Navier Stokes Equations in Radial Coordinates}
\subsection{Bibliography}
\end{document}
